\begin{frame}{$F(x)+x$ 의 차원 매칭 + 항등 블록의 시작점}
  \begin{block}{실무: 차원 매칭}
    \small
    더하려면 $F(x)$ 와 $x$ 의 차원(채널 $\times$ 공간)이
    \textbf{정확히 일치}해야 한다. ``채널은 늘고 공간은 줄어든다''
    패턴을 따르면 $F$ 가 64$\to$128 채널, 28$\times$28$\to$14$\times$14
    변환할 때 $x$ 는 그대로 --- \kb{1$\times$1 합성곱}(공간 유지,
    채널만 바꿈)으로 $x$ 의 차원을 맞춰준 뒤 더한다
    (필요하면 2$\times$2 스트라이드 1$\times$1 합성곱으로
    공간도 동시 축소).
  \end{block}
  \vskip0.5em
  \begin{block}{``항등 블록''이 왜 좋은 시작점인가}
    \small
    $F$ 를 \textbf{정규화된(0 근처) 작은 값}으로 초기화하면
    학습 \textbf{시작} 시점에 $F(x)\approx 0$ 이라
    $y\approx x$ --- \textbf{항등 함수로 시작}. ``깊은 네트워크가
    최적의 시작점을 얕은 네트워크의 항등 함수로 가질 수 있다''
    $\Rightarrow$ 학습이 ``0에서 $F(x)$ 를 만들어야 한다''가 아니라
    ``이미 잘 동작하는 항등 함수 위에서 \textbf{잔차만} 미세 조정''.
    ResNet이 ``깊은 망이 왜 어려운가''를 \textbf{표현력} 문제가
    아니라 \textbf{최적화(초기화) 문제}로 재정의한 것.
  \end{block}
\end{frame}
